\documentclass[11pt,a4paper]{article}
\usepackage[margin=1.8cm]{geometry}
\usepackage{amsmath,amssymb,array,tabularx,fontspec,cancel}
\setmainfont{FreeSerif}
\setlength{\parindent}{0pt}\setlength{\parskip}{0pt}
\renewcommand{\arraystretch}{1.15}
\newcounter{sq}
\newcommand{\sq}[1]{\stepcounter{sq}\par\smallskip\noindent\hangindent=11em\hangafter=1
  \makebox[11em][l]{\textbf{\Roman{sq}.}}#1\par}
\newcommand{\dtab}[2]{\par\smallskip\noindent\hspace*{11em}\begin{tabular}{#1}\hline #2\\ \hline\end{tabular}\par\smallskip\hspace*{11em}}
\begin{document}
\begin{center}
  {\Huge\bfseries The Hope Science Academy}\\[3pt]
  {\small Excellence is our standard not goal}\\[2pt]
  {\small 247-E-1, (Opposite LDA School) Johar Town, Lahore. PH: 0300-8194789, 0322-4157001}
\end{center}
\medskip\noindent
\begin{tabularx}{\linewidth}{|X|X|X|X|}
\hline
\textbf{Student Name:} & \textbf{Roll Num:} & \textbf{Class:} 11th & \textbf{Paper Code:} \\ \hline
\textbf{Subject Name:} Statistics & \textbf{Time Allowed:} & \textbf{Total Marks:} 1108 & \textbf{Exam Date:} 02-Nov-2026 \\ \hline
\multicolumn{4}{|l|}{\textbf{Exam Syllabus:} FULL BOOK} \\ \hline
\end{tabularx}
\bigskip
\noindent\textbf{Q2. Write short answers of the following questions.}\hfill\textbf{1$\times$668=668}

\sq{Define the word statistics and explain its different meanings?}
\sq{Define the following term: Population and sample}
\sq{Define the following term: Parameter and statistic}
\sq{Distinguish between qualitative and quantitative variables.}
\sq{Write the following using a summation sign with appropriate index: $Y_{3}+Y_{4}+\ldots+Y_{15}$}
\sq{Write the following using a summation sign with appropriate index: $Y_{1}^{2}+Y_{2}^{2}+Y_{3}^{2}+Y_{4}^{2}$}
\sq{Write the following using a summation sign with appropriate index: $(Y_{1}-\mu)^{2}+(Y_{2}-\mu)^{2}+(Y_{3}-\mu)^{2}$}
\sq{Write the following using a summation sign with appropriate index: $bY_{20}+bY_{21}+\ldots+bY_{30}$}
\sq{Expand the following summation and product sign: $\sum_{i=1}^{5}Y_{i}$}
\sq{Expand the following summation and product signs: $\sum_{i=5}^{8}(Y_{i}-\mu)$}
\sq{Expand the following summation and product signs: $\sum_{i=1}^{3}Y_{i}^{2}$}
\sq{Expand the following summation and product signs: $\sum_{i=2}^{4}(Y_{i}-9)$}
\sq{Expand the following summation and product signs: $\prod_{i=1}^{3}Y_{i}$}
\sq{Expand the following summation and product signs: $\prod_{i=3}^{6}(aY_{i})^{2}$}
\sq{Expand the following summation and production signs: $\prod_{i=2}^{4}(X_{i}-Y_{i})$}
\sq{Classify the following as categorical, discrete or continuous variable: Sex of an insect}
\sq{Classify the following as categorical, discrete or continuous variable: Weights of plants}
\sq{Classify the following as categorical, discrete or continuous variable: Major crops of Pakistan}
\sq{Classify the following as categorical, discrete or continuous variable: Level of satisfaction}
\sq{Classify the following as categorical, discrete or continuous variable: Teaching standards}
\sq{Classify the following as categorical, discrete or continuous variable: Temperature measured in Fahrenheit}
\sq{What are the meanings of statistics?}
\sq{Define data.}
\sq{What is datum?}
\sq{Define sampling?}
\sq{What is ratio?}
\sq{What is proportion?}
\sq{How parameters are represented?}
\sq{What is sampling variability?}
\sq{Define experiment.}
\sq{What is sample survey?}
\sq{What is order statistic?}
\sq{What is model?}
\sq{What is random error?}
\sq{What is sigma $(\sum)$ ?}
\sq{Write $Y_{5}+Y_{6}+Y_{7}+\ldots+Y_{10}$ using a summation notation with suitable index.}
\sq{Write $y_{1}^{2}+y_{2}^{2}+y_{3}^{2}+y_{4}^{2}+y_{5}^{2}$ using a summation notation with suitable index.}
\sq{Expand $\sum_{i=10}^{15}(y_{i}+3)$}
\sq{Expand $\sum_{i=1}^{8}y_{i}^{3}$}
\sq{What is product $(\pi)$ ?}
\sq{Write $y_{6}^{2}\times y_{7}^{2}\times y_{8}^{2}\times y_{9}^{2}\times y_{10}^{2}$ using product notation with suitable index.}
\sq{Write $(y_{4}-\mu)^{2}\times(y_{5}-\mu)^{2}\times(y_{6}-\mu)^{2}$ using product notation with suitable index.}
\sq{Expand $\prod_{i=1}^{8}by_{i}^{2}$}
\sq{Expand $\prod_{i=1}^{5}(x_{i}-y_{i})$}
\sq{Evaluate n!.}
\sq{Solve 7!.}
\sq{Solve 5!.}
\sq{Write down the importance of statistics in four disciplines.}
\sq{When a variable is fixed?}
\sq{When a variable is random?}
\sq{What are the types of variable?}
\sq{Define qantitative variable.}
\sq{Define qualitative variable.}
\sq{What are the types of quantitative variable?}
\sq{Define continuous variable.}
\sq{define discrete variable.}
\sq{Give some examples of continuous variable.}
\sq{Give some examples of discrete variable.}
\sq{What are the sources of data?}
\sq{What do you mean by statistical method?}
\sq{What is meant by primary data?}
\sq{Define inferential statistic.}
\sq{Explain discrete and continuous data}
\sq{Given $x_{i}$=5, 7, 10, 12 and $y_{i}$=8, 13, 18. Compute $\sum X_{i}^{2}$ and $\sum X_{i}Y_{i}$}
\sq{Define Statistics.}
\sq{Name the sources of secondary data.}
\sq{Give the uses of statistics. (Only two)}
\sq{Expand the following summation. (a) $\sum$ i = 15yi  (b) $\sum_{i=1}^{5}y_{i}$}
\sq{Differentiate between quantitative and qualitative data.}
\sq{Differentiate between population and sample.}
\sq{Differentiate between a variable and a random variable.}
\sq{Define statistics in singular sense.}
\sq{Define the term error. Name the types of error.}
\sq{Define secondary data with an example.}
\sq{What is the difference between constant and variable?}
\sq{Define statistic.}
\sq{Define qualitative variable.}
\sq{Distinguish between statistics and statistic.}
\sq{What are the two types of quantitative variables?}
\sq{What is description statistics?}
\sq{Define the term population.}
\sq{Enlist three main methods of collection of primary data.}
\sq{Write the names of two branches of statistics.}
\sq{Write two examples each for continuous and discrete variable.}
\sq{Define variable.}
\sq{Write two properties of G.M (Geometric Mean).}
\sq{What is sample?}
\sq{Define parameter.}
\sq{What is statistic?}
\sq{Define constant.}
\sq{Define variable.}
\sq{What is secondary data?}
\sq{Define statistics in plural sense.}
\sq{Expand the following $\sum_{i=2}^{4}(x_{i}-y_{i})$}
\sq{What are different methods of representation of statistical data?}
\sq{What is different method of representation of statistical data}
\sq{What is one-way classification?}
\sq{What is two-way classification?}
\sq{What is three-way classification?}
\sq{Define many way classification.}
\sq{What is tabulation?}
\sq{Define title.}
\sq{What are column captions?}
\sq{Define box-head.}
\sq{Define row captions.}
\sq{What is stub?}
\sq{What is body of table?}
\sq{What is prefatory note?}
\sq{What are footnotes?}
\sq{What is source note?}
\sq{Define frequency distribution.}
\sq{How a frequency distribution is constructed? i i}
\sq{What are open-end classes?}
\sq{What are class limits?}
\sq{What is class interval?}
\sq{Define class frequency.}
\sq{Define class mark.}
\sq{Define class boundary.}
\sq{What is cumulative frequency?}
\sq{Define relative frequency.}
\sq{Write down some common uses of graphs.}
\sq{Write down some important points for drawing graphs.}
\sq{What do you know about simple bar diagram?}
\sq{What is multiple bar diagram?}
\sq{What is component bar diagram?}
\sq{What is pie-diagram?}
\sq{Define histogram.}
\sq{Give some advantages of histogram.}
\sq{What is frequency polygon?}
\sq{How frequency polygon constructed?}
\sq{What is frequency curve?}
\sq{What is scatter plot?}
\sq{What is bivariate frequency distribution?}
\sq{In pie-diagram, how can we obtain angle for each sector?}
\sq{Write down the steps involved in the construction of simple bar diagram.}
\sq{Write down the steps involved in the construction of multiple bar diagram.}
\sq{Write down the steps involved in the construction of component bar diagram.}
\sq{Write down the steps involoved in the construction of pie-diagram.}
\sq{Write down the steps involved in the construction of histogram.}
\sq{Write down the steps involved in the construction of frequencies polygon.}
\sq{Write down the steps involved in the construction of ogive.}
\sq{Give the frequencies of 5 classes 2, 6, 7, 3, 4, find the cumulative frequencies and relative frequencies.}
\sq{What are the class boundries and size of class interval of (i) 3.5-4.6 (ii) -2- (+2)?}
\sq{Out line the four major parts of s table.}
\sq{Give the frequencies of 5 classes 3, 7, 8, 5 6, find cumulative frequencies, de-cumulative frequencies and relative frequencies.}
\sq{The observations about the number of rooten potatoes from twenty equal sized samples taken from a store are available as follows: 1,2,3,4,3,0,1,2,3,1,1,0,2,1,0,2,3,0,0,1,3. Make frequency table. Also find relative frequencies.}
\sq{How a frequency polygon is constructed?}
\sq{Enlist the main parts of a table.}
\sq{Differentiate between continous and discrete variable.}
\sq{Define classification.}
\sq{Define histogram.}
\sq{Define tabulation.}
\sq{What is simple bar chart?}
\sq{What are the main parts of a table? (Any two only.}
\sq{What is ogive?}
\sq{Define pie chart.}
\sq{Define relative frequency.}
\sq{Explain the term ``Open Ended'' classes.}
\sq{What is frequency distribution?}
\sq{Differentiate between histogram and historigram.}
\sq{Explain prefatory note and foot note in a table.}
\sq{What is class mark?}
\sq{Differentiate between ungrouped and grouped data.}
\sq{What are the important graphs of frequency distribution?}
\sq{What is size of class interval?}
\sq{What is frequency polygon?}
\sq{Differentiate between class mark and class width.}
\sq{Define Venn diagram.}
\sq{What are the class limits?}
\sq{What graph can be made from quantitative data, name any three graphs?}
\sq{Give three examples of tabular presentation.}
\sq{Define classification.}
\sq{Define histogram.}
\sq{Define tabulation.}
\sq{What is a simple bar chart?}
\sq{What are the main parts of a table? (only any two).}
\sq{What is ogive?}
\sq{Define pie chart.}
\sq{Define relative frequency.}
\sq{Explain the term ``Open Ended'' classes.}
\sq{What is frequency distribution.}
\sq{Explain prefatory note ans foot note in a table.}
\sq{What is a class mark?}
\sq{Differentiate between ungrouped and grouped data.}
\sq{What are the most graphs of frequency distribution?}
\sq{What is size of class interval?}
\sq{What is frequency polygon?}
\sq{Differentiate between class mark and class width.}
\sq{Define venn diagram.}
\sq{What are the class limits?}
\sq{What graph can be made from quantitative data, name any three graphs?}
\sq{Give three examples of tabular presentation.}
\sq{Find the geometric mean 50, 67, 39, 40, 36, 60, 54, 43.}
\sq{Define Average.}
\sq{What is measure of location?}
\sq{Give some properties of good average.}
\sq{Enlist the main measures of Central tendency.}
\sq{Write down the formula for population mean.}
\sq{Write down the formula for sample mean.}
\sq{Calculate mean for 87,91,89,88,89,91,87,92,90,98.}
\sq{Calculate mean for 10,15,20,25,30,35.}
\sq{For 2,4,5,6,7,6 show that $\sum(y-\bar{y})=0$.}
\sq{Prove that $\sum(y-\bar{y})=0$. (For upground data).}
\sq{Prove that $\sum f(y-\bar{y})=0$. (for grouped data).}
\sq{Show that mean is effected by the change of scale.}
\sq{Write down the coding formula for mean.}
\sq{Write down the short formula for mean.}
\sq{Give some merits of mean.}
\sq{Give some demerits of mean.}
\sq{What is weighted mean?}
\sq{How we calculate weighted mean.}
\sq{If $\sum wy=9326.6$ and $\sum w=132$, then cal [text cut off in the source]}
\sq{How we calculate geometric mean?}
\sq{Calculate G.M for 0.5, 10.0, 2.7, 3.48 and 4.7.}
\sq{Calculate the G.M for 45,30,335,40,44,332,42,37.}
\sq{Write down the properties of G.M.}
\sq{Give some merits of geometric mean.}
\sq{Give some demerits of a geometric mean.}
\sq{When geometric mean is appropriate?}
\sq{When harmonic mean is useful?}
\sq{How H.M is calculated?}
\sq{Calculate H.M for 10,20,30,40,50}
\sq{Calculate G.M for 3,8,11,0 and 10.}
\sq{Calculate H.M for 3,8,11,0 and 10.}
\sq{Write down the properties of harmonic mean?}
\sq{How we denote median?}
\sq{How we calculate median for odd data?}
\sq{Find $\sum y$ if $\bar{y}=5$, n=10.}
\sq{How can we calculate median for even data?}
\sq{Calculate the median for 87,88,89,89,90,91,91,92,98.}
\sq{How can we calculate median for grouped data.}
\sq{What are the merits of median?}
\sq{Define quantities.}
\sq{How can we calculate quantiles?}
\sq{What are quartiles?}
\sq{Define lower quartile.}
\sq{How we calculate lower quartile?}
\sq{How we define uppar quartile?}
\sq{How we calculate uppar quartile?}
\sq{Calculate $Q_{1}$ from 84.60,88.60,94.90,94.50,95.05.}
\sq{What are deciles?}
\sq{How we donate deciles?}
\sq{Where we calculate decile for grouped data?}
\sq{How we calculate percentiles for grouped data?}
\sq{How we find mode for ungrouped data?}
\sq{Write down the merits of mode.}
\sq{When a distribution is symmetrical?}
\sq{What is the empircial relation between mean.}
\sq{If mode=15, median =12, find them.}
\sq{If mean =45, median =30 find mode}
\sq{State some advantages of the mode.}
\sq{The sum of deviation of 10 values from 15 is 20. Find arithmetic mean.}
\sq{If mean =20, median=18.67, find mode.}
\sq{Find G.M if x=1,3,9.}
\sq{Given $\sum f=120$, $\sum fx=296$ and mode =2.94. Comment on the skewness}
\sq{If mean =20, median=18, find mode for a set of data?}
\sq{If a value of $Q_{2}$, $d_{5}$ and $P_{50}$ are equal to 72.32, then find the median of distribution}
\sq{Find the mode of 3,3,7,8,10,11,10,12.}
\sq{The sum of deviations of 15 values from 20 is 45. Find arithmetic mean.}
\sq{Write three types of Average.}
\sq{For a frequency distribution of a variable x, it is given that x=10+5u, $\sum fu=125$ $\sum fu=-45$, find the value of mean}
\sq{Find median of +1-3,3, -2 and 5.}
\sq{A distribution consists of 3 components with respective size 42,40 and 65 along with their means 2,2.5 and 5. Compute the combined mean.}
\sq{If mode =50, median=12, find mean.}
\sq{Give various methods of a computing arithmetic mean for grouped data.}
\sq{Why is difficult to calculate geometric mean for (a) 0,4,8 (b) 3,5,-2,1}
\sq{If $x_{1}$ =1 and $x_{2}$ =8 tjat A.M$>$G.M.}
\sq{Write down the empirical relation between mean, median and mode.}
\sq{What is the median of the word ``FREQUENCY''?}
\sq{In a single-peaked symmetrical distribution, mean =5, median =5 then what will be the value of mode?}
\sq{If mean =20, median=18.67, find mode.}
\sq{Find G.M x=1,3,9.}
\sq{Define G.M.}
\sq{Write two qualities of good average.}
\sq{Define mode with formula.}
\sq{Some of deviations of 25 values from 20 is 45. find A.M.}
\sq{If mode =15 and mean 10.5 find mediam.}
\sq{What is the relationship between mean, median and mode in a positively skewed distribution?}
\sq{Write down the formulas for different averages.}
\sq{Define the median.}
\sq{Write down empirical relationship b/w mean and mode}
\sq{Give $\ell=200$, $f_{m}$=25 $f_{1}$ 20, $f_{2}$=20 and h=10. find mode.}
\sq{Define harmonic mean.}
\sq{In a symmetrical mean is 35. Find the value of mediann and mode.}
\sq{Discuss advantages of median.}
\sq{Find mode the letter ``STATISTICS''.}
\sq{What is an array?}
\sq{Average of 5 values is 70. Find the sum of values.}
\sq{If mean is 20, Median =18. Find Mode.}
\sq{Give two properties of arithmetic mean.}
\sq{Observations. Find combined mean. [data missing in the source]}
\sq{Differentiate between simple arithmetic and weighted arithmetic mean.}
\sq{Calculate the geometric mean 1,1,1,27.}
\sq{Define Arithmetic mean.}
\sq{Calculate the geometric mean 1,1,1,27.}
\sq{Define Arithmetic mean.}
\sq{Compute the median for the data -2,5,0,-1,2,2.}
\sq{In single paked symmetrical distribution mean =5, median =5, then what will be the value of mode?}
\sq{A distribution consists of 3 components with respective size 45, 40 and 65, along with their means 2,202 and 2.5 and cmputer the combined mean.}
\sq{Write down merits of mode (any two)}
\sq{If $x_{1}$=2, $x_{2}$=8 show that A.M$>$G.M$>$H.M.}
\sq{The arithmetic mean of 100 items is 5 and the arithmetic mean of 150 item is 4.5 Find combined mean.}
\sq{What are two major properties of Median?}
\sq{In symmetrical distribution mode 20. Find value of Mean and Median.}
\sq{Give that median = 15 and mode = 20 by using the empricial relationship, find the value of mean.}
\sq{Define the quartiles and percentages.}
\sq{Find the mean and sum of the first n natural numbers.}
\sq{What do you mean by statistical average?}
\sq{Define the term measures of central tendency.}
\sq{Calculate median from the given data. 88.03,94.50,95.05,84.60.}
\sq{Define deciles and percentiles.}
\sq{In a moderately a symmetrical distribution the value of median is 42 and the value of mode is 40. Find the mean.}
\sq{Give $X_{1}$ 4, $X_{2}$=4, show that A.M = G.M = H.M.}
\sq{Given $\bar{x}_{1}=2.53,\ \bar{x}_{2}=1.97,\ \bar{x}_{3}=4.05,\ \bar{x}_{4}=$ [text cut off in the source] each mean is based on eight values}
\sq{Give the three dis-advantaged of H.M.}
\sq{Give Various methods of computing A.M.}
\sq{Find A.M when D=X-18 $\sum fD$ =150, n=25}
\sq{If A.M = 80, H.M =60, Find G.M.}
\sq{Given $\sum(x-10)$ =2.8 and n =5, find the sample mean.}
\sq{Given $X_{1}$ =4, $X_{2}$ =16. Show that G.M = $\sqrt{A.M\times H.M}$}
\sq{The sum of deviation of 10 values from x=40 is 250, what is the value of arithmetic mean?}
\sq{Compute geometric mean of 5,25,125.}
\sq{If X1 2= and X2 =8, then shpoe that A.M$>$G.M$>$H.M.}
\sq{If mean =75 and mode =70 find median.}
\sq{State when it is impossible to calculate Harmonic Mean.}
\sq{The Arithmetic mean of 20 values is 40.5 what is the total of values?}
\sq{Calculate geometric mean of X=1,1,27.}
\sq{Write down any two demerits of geometric mean.}
\sq{Write down two merits and two de-merits of Harmonic Mean.}
\sq{In a moderately skewed distribution, Mean =25 mode=31, Find value of Median.}
\sq{For a frequency distribution of a variable X, iit is given that X=10+5$\mu$, $\sum f=125$, $\sum fu=45$. Find te value of mean}
\sq{Write down two demerits of A.M.}
\sq{If G.M two values is 3. Find the product of two values.}
\sq{Given data $\sum fu$, $=1$, $\sum f=30$ and $u$, $\frac{Xi-98}{5}$ F arithmetic mean of 'X'.}
\sq{What are demerites of arithmetic mean?}
\sq{In moderately skewed distribution, mode=15, median=12, Find its mean.}
\sq{Why are Average called Measures of Central Tendency?}
\sq{Write down the name of Average.}
\sq{What are the types of dispersion?}
\sq{Wha is Absolute dispersion?}
\sq{What is relative dispersion?}
\sq{How can we calculate relative dispersion?}
\sq{What are the different types of absolute dispersion?}
\sq{Define range.}
\sq{How we calculate range from grouped data?}
\sq{How we calculate range from grouped data?}
\sq{What are the merits of range?}
\sq{What are the demeirts of range?}
\sq{How can calculate coefficient of range?}
\sq{Calculate range for the following observations: 13,3,7,15,17,5,23,27.}
\sq{What is quartile deviation?}
\sq{How we calculate inter-quartile range?}
\sq{How we calculate inter-quartile range?}
\sq{If $Q_{1}$=88.03 and $Q_{3}$=94.90, find Q.D and I.Q.R}
\sq{Calculate quartile deviation for $Q_{1}$=92.3750, $Q_{3}$=102.58333}
\sq{How we calculated M.D for grouped data?}
\sq{How can we calculating M.D for grouped data?}
\sq{Write down the properties of mean deviation.}
\sq{What are the merits of mean deviation?}
\sq{What are the demirates of mean deviation?}
\sq{How can calculate coefficient of mean deviation?}
\sq{If M.D from Median =3.92 and $\bar{y}$=16.25, find coefficient of mean deviation.}
\sq{What is variance?}
\sq{How we calculate veriance for ungrouped data?}
\sq{How we calculate variance for grouped data?}
\sq{Define standard deviation.}
\sq{How standard deviation is donated?}
\sq{How can calculate standard deviation for ungrouped data?}
\sq{Give some demerits of standard deviation.}
\sq{Write down some properties of variance.}
\sq{What are moments?}
\sq{What are shappard's corrections?}
\sq{What is coefficient of variation? What are its uses?}
\sq{Define skewness.}
\sq{Define absolute dispersion.}
\sq{Given $\sum x=180$, $\sum x^{2}=6660$ and n=5. Find Co-efficient of variation.}
\sq{Write down the types of dispersion.}
\sq{Define range with formula.}
\sq{What do you mean by co-efficient of variation, write the formula also.}
\sq{Write the formula of Kari Pearson's coefficient of skewness.}
\sq{Define the term ``Moments''.}
\sq{If $Q_{1}$ = 88.03 and $Q_{3}$ 94.50. Compute coefficient of Q.D.}
\sq{Define moment rations. In which units are they expressed?}
\sq{The S.D of a symmetrical distribution is \%. Find its four moments about mean. [value missing in the source]}
\sq{If mean =25, CV =64\% what will be the variance?}
\sq{Find Bowley's coefficient of skewness. Given $Q_{1}$ =53.67, Median= 65.0, $Q_{3}$ =76.33.}
\sq{If $m_{2}$=5, $m_{3}$=-12.8. Find $b_{1}$ and comment}
\sq{Define relative dispersion.}
\sq{Compute the correct coefficient of variation when $\bar{X}$=67.45, $S_{2}$=8.5275 and class interval is 3}
\sq{Given $\bar{X}$=100, Mode =95 and S=10, find coefficient of skewness.}
\sq{The mean of a distribution is 3 while it second moment about zero is 24, find its standard deviation.}
\sq{Given $\sum f=120$, $\sum fx$=296 and mode=2.94. DIscuss the skewness}
\sq{Given $\sum f=100$, $\sum f|x-\bar{x}|=800$ and $\sum fx$ =1600, Find coefficient of mean deviation.}
\sq{Define kurtosis}
\sq{The first moment about x=15 is 2.5, the second moment about mean is 16. What is the value of mean and standard deviation?}
\sq{Give the formula for calculating combined standard deviation for three subgroups}
\sq{Give $\bar{x}$ =55, $S^{2}$ = 121. Find coefficient of variation}
\sq{if $Q_{1}$ = 88.06, $Q_{3}$ =94.50, compute coefficient of Q.D}
\sq{A series comprises of hundred values each equal to ``5''. What will be the average and dispersion of the series?}
\sq{Describe the properties of variance (any two).}
\sq{What is symmetrical distribution?}
\sq{Define combined variance.}
\sq{Define standard deviation.}
\sq{What can you say about skewness if mean =1393 and mode =145.8.}
\sq{Write formula of Bowely's coefficient of skewness.}
\sq{Find variance of given data 5,5,5,5,5,5 and 5.}
\sq{If var(x) =8, find var (x+3).}
\sq{If $Q_{1}$ =136.62, $Q_{2}$ =146 and $Q_{3}$ =153.13 find bowley coefficient of skewness.}
\sq{If var(x)=8, find var(2x).}
\sq{Define mean deviation.}
\sq{If $Q_{1}$ =136.62, $Q_{3}$ =153.13, find co-efficient of quartile deviation.}
\sq{If S.D(x)=10 then find S.D (2x+10)=10?}
\sq{Define Dispersion.}
\sq{Given $Q_{1}$ =125 $Q_{3}$ =175. Find the coefficient of quartile deviation.}
\sq{When the dist. Is lepto kurtic and platy kurtic.}
\sq{When the dist. is lepto kurtic and platy kurtic.}
\sq{Compute range of the data -1,-3, -4, -10, -20.}
\sq{If some of squares of deviation from mean is 192 for 12 values, then find standard deviation of the values.}
\sq{The coefficient of variation of a distribution is 45\% find S.D if its mean is 16.}
\sq{The first two moments about a value x=1.5 of variable are 01 and 1.7. Calculate co-efficient of variation.}
\sq{In the first moment about the number 2 is equal to 5. What is the mean?}
\sq{The lower and upper quartile of a distribution are 142.36 and 167.73 respectively, while median is 153.50. Find co-efficient of skweness.}
\sq{What would the shape and name of the frequency distribution if (a) Mean=Median=Mode (b) Mean$>$Median$>$Mode.}
\sq{Define the semi inter quartile range.}
\sq{If lower quartile is 20 and quartile deviation is 5, Find upper quartile.}
\sq{Differentiate between positive and negative skewness.}
\sq{How is variance defined? Also give its formula.}
\sq{Define skewed distribution.}
\sq{What can you say about skewness if median=76.29 and two quartiles are 71.38 and 81.25.}
\sq{Give names of any four methods for calculating dispersion.}
\sq{Write any two properties of standard deviation.}
\sq{Find the range of -1, -3, 0, 2, 3,}
\sq{If var(X) -25, find var (2X+4).}
\sq{Define quartile deviation.}
\sq{if n=10, $\sum x$ =40 and S=2 then find coefficient of variation.}
\sq{If $Q^{1}$=10, $Q_{2}$=20 and $Q_{3}$=30, find the coefficient of skewness.}
\sq{What are two merits and one demirats of mean deviation?}
\sq{If a distribuion has mean 14.3 and mode 41487. What can you say about the skewness?}
\sq{If $Q_{1}$-20, Quartle Deviation =30, FInd $Q_{3}$.}
\sq{Given mean =50, Median =48, SD=6 Find coefficient of skewness.}
\sq{What do you understand by the term ``Quantiles''?}
\sq{Find quartile deviation and co efficient of quartile deviation of 7.4, 7.4, 7.4, 7.4, and 7.4.}
\sq{Define negtively skewed data. Give one example from real life.}
\sq{Find Bowley's coefficient of skewness if $Q_{1}$=84, $Q_{3}$=79 and median=81.}
\sq{In a modertely skewed distribution, mean=100, mode=95 and S.D=10. Find coefficient of skewness.}
\sq{Find K for the probality distribution.\dtab{|l@{\hspace{3em}}c@{\hspace{3em}}c@{\hspace{3em}}c|}{x & 0 & 1 & 2 \\ \hline p(x) & 3k & 2k & k}}
\sq{S.D of a distribution is 4,find 2nd moment about mean.}
\sq{If var (x) =10 and y=5x+20, then find var(y).}
\sq{Enlist various relative measures of dispersion.}
\sq{Can mean, median and mode be same, if yes, state in what situation?}
\sq{If first three moments about X=20 of a distribution are: 1, 4, 40, then find the value of ``b1''.}
\sq{What is the range of Bowley's coefficient of skewness?}
\sq{If $Q_{1}$=13.73, $Q_{3}$=38.29, compute Quartile deviation.}
\sq{Calculate range of 13, 3, 4, 15, 17, 5, 13, 27.}
\sq{What are different measures of absolute dispersion?}
\sq{If var(x) =4, Find var (3x).}
\sq{What is skewness?}
\sq{When distribution is symmetrical?}
\sq{When distribution is negatively skewed?}
\sq{When distribution is normal?}
\sq{When distribution is platykutic?}
\sq{When distribution is lepokurtic?}
\sq{For what values S.D and variance are equal?}
\sq{Give the Bowley's formula of a coefficient or skewness.}
\sq{What would be the shape of the distribution if Mean$<$Median$<$Mode?}
\sq{Which absolute measure of dispersion is the best?}
\sq{Define inter-quartile range and quartile deviation.}
\sq{Write four relative measures of dispersion.}
\sq{If coefficient of skewness is zero, then what will you say about mean, median and mode?}
\sq{Range of 5 observation is -6 and smallest observations is -2, find the largest observation.}
\sq{Given E(x)= 204 and C.V(x)=6, find S.D (x).}
\sq{Define Price Relatives.}
\sq{$\sum p_{1}q_{o}=50$, $\sum p_{o}l_{o}=425$. Fin index number}
\sq{Given most suitable average for index number, with reason.}
\sq{Given W= 20, 30, 40, 50, 60, l = 100, 105, 110, 105, 120. Find weighted index.}
\sq{What are quantity index numbers?}
\sq{Define composite index number.}
\sq{Define link relative.}
\sq{Write the Formula for Fisher's index no}
\sq{Given $\sum p_{o}l_{o}=1500$, $\sum p_{n}q_{o}=2040$. Find base year weighted index number.}
\sq{What are types of Composite index number?}
\sq{Write down some mein steps involved in construction of index number}
\sq{What is variable.}
\sq{What are random numbers?}
\sq{How we generate random numbers?}
\sq{What is pseudo random technique?}
\sq{What is the application of random numbers?}
\sq{Give some examples of discretes random variable.}
\sq{Define continuous random variable.}
\sq{Give some examples of continuous random variable.}
\sq{What are the types of random variable?}
\sq{Which methods are used in drawing random numbers?}
\sq{Define random experiment.}
\sq{Define continuous random variable?}
\sq{What is mean by random variable?}
\sq{Can a random variable assume negative values?}
\sq{What is the application of random numbe?}
\sq{What is discrete random variable?}
\sq{Make random variable for tail when three coins are tossed.}
\sq{Write down the assumptions of random numbers.}
\sq{Differentiate between a variable and random variable.}
\sq{Write down the properties of a continuous random vireable.}
\sq{How random number can be generated?}
\sq{Give two examples of random variable.}
\sq{Describe the two different types of random variable.}
\sq{What is the same sample space of a random experiment?}
\sq{Properties of discrete random variable.}
\sq{Make a random variable for number of heads when two coins are tossed.}
\sq{Two dice are thrown make a random variable for sum of dots.}
\sq{Define random numbers.}
\sq{What are the expectations and standard deviation of a constant?}
\sq{Given a random variable X with E (x) = 12 and S.D (X) =5. Find E(x2).}
\sq{What does P.D.F stands for?}
\sq{Write three properties of random experiment.}
\sq{Define discrete random variable.}
\sq{What is discrete random variable?}
\sq{Give some examples of dicrete random variable?}
\sq{Define continuous random variable.}
\sq{Give some examples of countinuous random variable.}
\sq{How does we describe probability mass function?}
\sq{How we calculated variance of variable y?}
\sq{If E(x)=1.1 and E($X^{2}$) =2.1, find Var (x)}
\sq{If E(x) =1.1, Find E(3x+5).}
\sq{\dtab{|l@{\hspace{3em}}c@{\hspace{3em}}c@{\hspace{3em}}c|}{X & 30 & 35 & 40 \\ \hline P(X) & 15/45 & 24/45 & 6/45} Find E(x)}
\sq{If E(x)=7, E($x^{2}$)=54.83, find Var (x)}
\sq{If E(x)=200, C.V(x)=7\%, find Var (x)}
\sq{If E(y)=2.1, E($y^{2}$)= 9.8, find Var (y)}
\sq{If Var (x)=2.07, find 2.07 S.D (x).}
\sq{If E(x) = 9.94, S.D (x)=1.44, find C.V (x).}
\sq{Head: Tail=3:1, find C.V (x).}
\sq{If E(x) =5, find E(-3x+2).}
\sq{If E(5x+10)=18, $E(5x+10)^{2}$=530, find VAr (5x+10)}
\sq{Find E(5-3y), When E(y)=1.4.}
\sq{Check for Y=1,2,3,4. Is f(y) =$\frac{y+1}{14}$ a p.d.f?}
\sq{If Var (x)=5, Var (Y)=$\frac{1}{3}$, then find variance and S.D of X?}
\sq{x:0 1 2 P(x) : 4c 3c c find the value of c.}
\sq{If E(x)=1.1 and E($x^{2}$)=2.1, find Var (x)}
\sq{What is methametical expectations?}
\sq{Give the properties of probabilties distribution.}
\sq{If E(x) =2 and E($x^{2}$) =20. Find standard deviation of X}
\sq{If E(x) =3 and E(y) =2.5, find E(x+y).}
\sq{What is distribution function?}
\sq{Define probability destiny function (p.d.f).}
\sq{Write down properties of expectation, any two.}
\sq{Given E($x^{2}$)=400 and Var(x) =144, then find mean of X values (i.e.E(x)).}
\sq{Write down properties of probability density function.}
\sq{$S^{2}$ =3 and E(x) =3. find E($x^{2}$).}
\sq{Var (x)=2 and Var (y) =5. find var(3y-y) if X and Y are independent.}
\sq{What are the properties of discrete probability distribution?}
\sq{If E(x) =5, E(-3x+2).}
\sq{What is discrete probability distribution?}
\sq{Given E(x) =10 and $E(x)^{2}$ =500, find coefficient of variation.}
\sq{What are properties of distribution function?}
\sq{Given E(x)=0 and E(y) = $\frac{8}{9}$, find E(3x-2y).}
\sq{Given f(x)= $\frac{k}{x}$, x=1,2,3. find K.}
\sq{if E($x^{2}$) = 400, S,D(x)=12, find E(x)?}
\sq{if E(x) =0.63, Var (x) =0.2331, find E($x^{2}$)?}
\sq{if E(x) =0.63, Var (x) =0.2331, find E($x^{2}$)?}
\sq{If E(5x+10)=18, $E(5x+10)^{2}$ = 530, find var (5x+10)}
\sq{Give E(x) =2.5, E($x^{2}$)=7.5, find variance and coefficient of variation}
\sq{Give E(x) =2.5, E($x^{2}$)=7.5, find variance and coefficient of variation}
\sq{Define countinuous properties of expectation.}
\sq{Define probability distribution along with its properties}
\sq{What is the mathematical expectation to win Rs. 50 if a head appears and loss Rs. 50 if a tail appears in a single toss of a fair coin?}
\sq{Given E(x)=0.55. Var (x)=1.35 and y=2x+1. Find E(y) and Var (y).}
\sq{if E(x) =3, E(y) =2.5 find E(2x=y)}
\sq{What is an expected value of a random variable?}
\sq{Given x=-10, -20, 30 and P(x)=2k, 3k, 5k, find K}
\sq{Given P(x)=C$|x-2|$ and x=-1,0, 1, 3 find E(x)}
\sq{Given x=0,2,3 and P(x)= $\frac{|1-x|}{4}$ find E(x).}
\sq{If E(X)=1.1, Find E(3x+5).}
\sq{Given x=0, 1, 2 and P(x) = $\frac{5}{8},\frac{4}{8},\frac{1}{8}$ is this a probability function?}
\sq{A countinuous random variable x has a density function f(x) = $\frac{cx}{4}$ for $1\leq x\leq 4$.}
\sq{If E(x)=1.15= then E(3X+5).}
\sq{If E(x) =1.1, E($x^{2}$)=2.1 find VAr (x)?}
\sq{If E(x) = 1.1, Find E(3x9+5)?}
\sq{Find E(x) when var (x)=4, E($x^{2}$) =20}
\sq{Given E($X^{2}$) = 400. S.D (x) =12, then find E(x).}
\sq{Given the E($X^{2}$) =400, S,D (X) =12 find E(X)}
\sq{What is probability mass function?}
\sq{Write down the properties of probability mass function.}
\sq{Find E(x), E($X^{2}$) from the following data [data missing in the source]}
\sq{A continuous random variable X has a density function f(x)= $\frac{Cx}{4}$ for a$<$x$<$4. Find the value of C.}
\sq{What does p.d.f stands for?}
\sq{Given E(x) =0 and E($x^{2}$)= $\frac{8}{9}$, Find E($2x^{2}-3x+5$)}
\sq{E(x)=1.1 and E($x^{2}$)=2.1, find Var (x)}
\sq{A die is rolled if x= face value. Find E(x).}
\sq{What will be the K in the given probability distribution? [table missing in the source]}
\sq{Given f(x)= $\frac{k}{x}$; x=1,2,3 find k}
\sq{If E(x)=10, then find E(5x-25).}
\sq{Given x=0,1,2 : P(x) = 4c, 3c, c, find the value of c}
\sq{Define trial.}
\sq{What is Binomial experiment?}
\sq{What are the properties of binomial properties?}
\sq{What are the parameters of the binomial distribution?}
\sq{What are binomial coefficient?}
\sq{When binomial distribution is skewed?}
\sq{When binomial distribution is positively skewed?}
\sq{When binomial distribution is negatively skewed?}
\sq{If n=4, p-$\frac{1}{2}$, find P(x=3).}
\sq{if you toss for a fair die 6 times, what is the probability of getting no even number?}
\sq{Out of 800 families with 5 children each. How many would you expect to have 4 boys?}
\sq{Give two properties of a Binomial experiment.}
\sq{In binomial distribution n=10 and p=0.5 find its variance.}
\sq{If n=10, n=4, k=6. Find variance of the hypergeometric distribution .OR find E(x) and S.D(x).}
\sq{If n=5, 0.25, Find P($x\leq 1$).}
\sq{Define Bernoulli trials.}
\sq{Criticize the statment that in a binomial distribution Mean =6 and standard deviation is 2.5}
\sq{Find mean and variance of following binomial distribution. $\left(\frac{1}{2}+\frac{2}{3}\right)^{2}$}
\sq{Define Hyper-geometric distribution.}
\sq{In a binomial distribution n=5 and p= 0.5 Find coefficient of variation .}
\sq{In a binomial distribution, mean=6 and variance=4. Find its parameter.}
\sq{Whar is the variance of Hypergeometric distribution?}
\sq{Given N=6, n=3 and K=4, find P(x=0).}
\sq{What is bimetrical experiment?}
\sq{What are the perameters of hyper-geometric distribution?}
\sq{If mean=20, P=0.2 in binomial distribution, then find n.}
\sq{What is mean and variance of biomial distribution?}
\sq{If n=8 and p= 0.4, then find variance of binomial distribution.}
\sq{if N=10, n=3, k=4, find variance of Hypergeometric distribution.}
\sq{What is binomial probability distribution ?}
\sq{What is the range of binomial r.v ``X''?}
\sq{If a fair coin is tossed 5 times, find the probability of getting at the most 2 heads?}
\sq{Define hyper-geometric experiment.}
\sq{N=10, n=2 and k=3. Find P(x=0).}
\sq{If n=8, p= $\frac{1}{3}$ in binomial probability distribution, then find mean and variance}
\sq{Define the perameters of the binomial probability distribution.}
\sq{Find the number of trials of a binomial distribution which has mean=12 and standard deviation=2.}
\sq{Define binomial Random variable.}
\sq{If N= 10 N-K=4 n=5., Find variance of Hypergeometric distribution.}
\sq{If E(x)=3 and Var (x)=2, then find parameters of binomial distribution of random variable X.}
\sq{If n= 12 and p= 0.6, then find mean and variance of binomial distribution.}
\sq{Is it possible to have binomial distribution with mean 10 and standard deviation 4?}
\sq{What is hyper-geometric random variable?}
\sq{If N=10, N-K=4, n=5 find P(X=1).}
\sq{If n=10 and P= 0.4, find mean and variance of binomial distribution.}
\sq{A random variable X has a binomial distribution with n=5 and p= 0.2 find p(x=2).}
\sq{If E(x)=2.4 and p=0.3, find the S.D of binomial distribution.}
\sq{Given N=8, n=3 and K=4. Find P(x=2).}
\sq{If N=10 and q=0.4, then find P(x=2).}
\sq{Is it possible to have a binomial distribution with mean = 5 and S.D =3.7?}
\sq{If X is a binomially distributed with n=10 and p=0.4, then find variance of Y= $\frac{x-10}{6}$}
\sq{For a binomial distribution with n=10 and p=0.5.. Find probability of successes.}
\sq{Given N=10, n=4 and k=3. find P(x=1).}
\sq{In a binomial distribution mean=2.4 and S.D=1.2. Find the value of ``n''.}
\sq{If x $\sim$ h (x:11,5,4). Find mean and variance of hyper- geometric distribution.}
\sq{If ``x'' is a binomial random variable with n=9, p= $\frac{1}{3}$ then find E(y) and Var(y).}
\sq{In a binomial distribution with n=5, what is the value of 'P' IfP(x=0) =P(X=1).}
\sq{If N=40, n=5, K=4, then find values of mean and variance of Hyper- geometric distribution.}
\sq{In a binomial distribution mean =38 and q=0.83, find n and p.}
\sq{Write down farmula for mean and variance of hyper-geomatric distribution.}
\sq{What will be the mean and variable of binomial distribution if n=6 and p= 0.6?}
\sq{If N=10, n=5, K=3 find mean of hypergeometric distribution by using formula of mean.}
\sq{State probability function of Hypergeometric distribution. Also write its mean and variance.}
\sq{Given N=10, n=2 and K=2. Find P(X=0)}
\sq{If n=4, p= $\frac{1}{2}$ find P(x=3)?}
\sq{What is binomial distribution?}
\sq{What is mean, variance and S.D of $(q+p)^{4}$?}
\sq{Given P(x) =k. $^{4}C_{x}$ and x= and x=0,1, 2, 3,4, find the values of K.}
\sq{In hypergeometric, N=40, n=5, K=4. Find mean and variance.}
\sq{Find the mean of hypergeometric distribution when N=8, n=2 and find K=8.}
\sq{If n=10 and p=0.5. Find mean S.D of binomial distribution}
\sq{In binomial distribution, if n=10 and p= $\frac{1}{4}$, then find variance of binomial distribution}
\sq{In binomial distribution, n=7, p=.0.7. find coefficient of variation.}
\sq{Under which circumstance we can apply the binomial distribution and the hypergeometric distribution?}
\sq{Give any two properties of hypergeometric experiment which are diferent from binomial experiment.}
\sq{If in binomial distribution Mean =3.2 Var=1.152. Find n and p.}
\sq{The sum of 10 items is 12 and sum of square equal to 26.9, then find standard deviation.}

\end{document}
